\documentclass[a4paper,12pt]{article}

\usepackage[french]{babel}
\usepackage[T1]{fontenc}
\usepackage[utf8]{inputenc}
\usepackage{amsmath, amssymb}
\usepackage{geometry}
\usepackage{array}

\geometry{top=2cm, bottom=2cm, left=2cm, right=2cm}

\begin{document}

\begin{center}\Large\textbf{CORRIGÉ DU SUJET 1}\end{center}

\section*{PREMIÈRE PARTIE : AUTOMATISMES -- QCM}

\subsection*{Question 1}
Si $a = 0,5$ alors $a^2 = 0,25$ et on a donc $a^2 \leq a$. \\
Si $a = 2$ alors $a^2 = 4$ et on a donc $a^2 \geq a$. \\
On observe donc que $a$ et $a^2$ ne sont pas toujours classés dans le même ordre.

\boxed{\text{Réponse D}}

\emph{Rappel : si $a > 0$ alors $a^2 < a$ si et seulement si $0 < a < 1$.}

\subsection*{Question 2}
\[
\frac{12}{30} \times 100 = \frac{2 \times 6}{5 \times 6} \times 100 = \frac{2}{5} \times 100 = 2 \times 20 = 40
\]

\boxed{\text{Réponse D}}

\subsection*{Question 3}
\[
(1,11 - 1) \times 100 = 11
\]

\boxed{\text{Réponse B}}

\subsection*{Question 4}
Sur le graphique, on lit que la courbe passe par le point de coordonnées $(0;1)$. Donc $f(0) = 1$.

\boxed{\text{Réponse B}}

\subsection*{Question 5}
La hauteur de la barre correspond aux effectifs. C'est donc l'anglais qui est la LV1 la plus suivie dans cet établissement.

\boxed{\text{Réponse C}}

\subsection*{Question 6}
« N'obtenir aucun 6 » est l'événement contraire de « obtenir au moins un 6 ». Sa probabilité est donc $1 - 0,19 = 0,81$.

\boxed{\text{Réponse A}}

\subsection*{Question 7}
Le choix se faisant au hasard, il y a équiprobabilité.
\[
P_{A_3}(L) = \frac{70}{80} = \frac{7}{8}
\]

\boxed{\text{Réponse C}}

\subsection*{Question 8}
$350$ km/h est l'ordre de grandeur de la vitesse maximale d'un train à grande vitesse circulant avec des passagers.

\boxed{\text{Réponse C}}

\subsection*{Question 9}
$f$ est une fonction affine non constante. Le coefficient directeur vaut 3 donc $f$ est strictement croissante. $f(x) = 3x - 1$ s'annule en $x = \dfrac{1}{3}$.

\boxed{\text{Réponse A}}

\subsection*{Question 10}
\[
\frac{2}{7} \times s = 30 \Rightarrow s = \frac{7}{2} \times 30 = 105
\]

\boxed{\text{Réponse B}}

\subsection*{Question 11}
Coefficients multiplicateurs : $0,7$ puis $0,8$. $c = 0,7 \times 0,8 = 0,56$. $t = (0,56 - 1) \times 100 = -44\%$.

\boxed{\text{Réponse B}}

\subsection*{Question 12}
$\dfrac{6}{3} = 2$ donc $N(3;3)$ n'appartient pas à la courbe.

\boxed{\text{Réponse B}}

\newpage

\section*{DEUXIÈME PARTIE}

\subsection*{Exercice 1}

\begin{enumerate}
\item Lecture graphique :
\begin{enumerate}
\item[a)] $h_0 = 2$ mètres.
\item[b)] $h_m = 18$ mètres.
\item[c)] $t_m = 2$ secondes.
\item[d)] $t_1 \approx 4,1$ secondes.
\end{enumerate}

\item Par le calcul :
\begin{enumerate}
\item[a)] $h(0) = -4 \times 0^2 + 16 \times 0 + 2 = 2$.
\item[b)] Maximum atteint en $t = -\dfrac{b}{2a} = -\dfrac{16}{2 \times (-4)} = 2$ secondes. $h(2) = -16 + 32 + 2 = 18$ mètres.
\item[c)] $h(t) = 0 \Rightarrow -4t^2 + 16t + 2 = 0$. $\Delta = 256 + 32 = 288$. $t_1 = \dfrac{-16 - \sqrt{288}}{-8} = \dfrac{16 + 12\sqrt{2}}{8} = 2 + 1,5\sqrt{2} \approx 4,12$ secondes.
\end{enumerate}
\end{enumerate}

\subsection*{Exercice 2}

\subsubsection*{Partie A}

\begin{enumerate}
\item $a_0 = 1024$, $a_1 = 512$, $a_2 = 256$.
\item À chaque pliage, l'aire est divisée par 2 : $a_{n+1} = \dfrac{1}{2} a_n$.
\item Suite géométrique : $a_n = 1024 \times \left(\dfrac{1}{2}\right)^n = \dfrac{1024}{2^n}$.
\item $a_{10} = \dfrac{1024}{2^{10}} = \dfrac{1024}{1024} = 1$ cm$^2$.
\end{enumerate}

\subsubsection*{Partie B}

\begin{enumerate}
\item $S_n = a_0 \times \dfrac{1 - \left(\frac{1}{2}\right)^{n+1}}{1 - \frac{1}{2}} = 1024 \times \dfrac{1 - \frac{1}{2^{n+1}}}{\frac{1}{2}} = 2048 \left(1 - \frac{1}{2^{n+1}}\right)$.
\item D'après le tableau, $S_n$ se rapproche de $2048$ lorsque $n$ augmente. On conjecture que $S_n$ tend vers $2048$ quand $n$ tend vers $+\infty$.

Interprétation : La somme des aires des carrés successifs tend vers le double de l'aire initiale. Cela correspond à la somme des aires de tous les pliages possibles.
\end{enumerate}

\end{document}